\section{宪制治理与协同演化 (ARI-RQGM)}\label{sec:rqgm}

此前各章都将 ARI 的\term{制度}——由提示词定义的、提出研究方向、
评审节点并评判质量的那些角色——视为在一次运行的整个生命周期内固定
不变。本章描述 \term{ARI-RQGM}:一种可选启用的执行模式,它允许这套
制度在一部自身不演化的宪制之下演化。该模式将 \cref{sec:exploration}
中原封未动的最佳优先树搜索包裹进一个基于纪元的治理层
(\texttt{ari.rqgm} 包):提示词与组件相互竞争、攻击、辩护,并被采纳
或退役——但每一次制度状态变更都要由一个确定性的、不调用 LLM 的
内核,对照冻结在代码中的规则表加以校验。默认模式仍是无治理的循环;
治理保持未配置时,不会构造任何治理对象,检查点与前述各章所描述的
逐字节相同。启用与否是一项运行级决定,在运行开始之前作出,并记录
于检查点。

\subsection{动机与源流}\label{sec:rqgm:motivation}

制度角色本身就是由提示词定义的产物,因此让它们演化,有望获得与搜索
循环从实验中萃取的同一类改进。然而,不受约束的自我修改恰恰是哥德尔机
(G\"odel machine)文献所警示的:能够改写自身评估器的组件,终将给
自己发放奖励。红皇后哥德尔机 \cite{iacob2026red} 加入了协同演化的
制衡——智能体之所以进步,正\emph{因为}对抗者与评估器在与之共同
演化——ARI-RQGM 的名称与核心动力学都承袭自这一工作脉络。但仅有协同
演化并不能约束对研究系统真正要紧的失效模式:分数虚高、指标操纵、
证据伪造,以及角色之间的合谋。

因此,ARI-RQGM 将协同演化循环置于一个永不演化的固定宪制层之下。
贯穿始终的设计立场是\emph{只阻断制度,绝不阻断研究}:治理可以否决
其自身的状态变更——采纳、处分、注册表写入——但绝不否决任何节点的
执行,而且搜索路径上的每个治理钩子都是故障开放(fail-open)的,
因此治理失效降级的是治理,而非实验。这是
\cref{sec:overview:principles} 的确定性原则向制度变更的延伸:内核
裁定、迁移解析与前沿修复都是纯函数——没有随机性,没有依赖挂钟的
决定——因此受治理的运行会重放出同一部制度史,正如无治理的运行会
重放出同一棵树。

\subsection{三个层级,四个门面}\label{sec:rqgm:layers}

治理将系统分层为三个层级:改变事物的权限严格递减,而被改变的自由
严格递增。

\begin{description}
  \item[\textbf{第 0 层,宪制层(固定)。}] 确定性内核、固定的结果
        核验路径,以及只追加的审计日志。它们从不演化;其执行的规则
        表存于代码之中,由一个宪制哈希(\cref{sec:rqgm:constitution})
        钉住,并仅出于溯源目的登记于组件注册表。
  \item[\textbf{第 1 层,制度层。}] 可演化的服务角色:生成器、评审者、
        对抗者、辩护者、裁决者与路由器。这些角色的提示词可以经由
        \cref{sec:rqgm:prompts} 的生命周期被替换,且仅能在纪元边界
        处替换。
  \item[\textbf{第 2 层,元层。}] 演化第 1 层的智能体——提示词
        变异器与净室生成器,外加更窄的辅助角色——它们自身也受治理,
        且权限严格小于其所修改的层级:元智能体可以提议一个替代者,
        却永远不能任命它。
\end{description}

四个门面掌管这套机制。\term{宪制内核}校验记录、哈希、能力、迁移、
角色分离、擦除、审计日志完整性、净室输入与上下文作用域——十二个
封闭的入口点,无一调用模型或读取时钟。\term{治理编排器}运行纪元
边界审计(\cref{sec:rqgm:boundary})并产出一份建议性报告;它从不
写入注册表。\term{注册表迁移引擎}是组件与提示词状态的唯一写入者,
对照一张固定迁移表解析变更。\term{前沿修复引擎}在发生退役后恢复
搜索前沿(\cref{sec:rqgm:boundary})。环绕它们的是各项逐节点服务:
提案路由器、对抗轮、提示词演化流水线、净室协调器,以及一个预算
管理器——后者可以让任何治理决策点降级,却绝不触碰节点执行。

\subsection{启动:模式解析与创始注册}\label{sec:rqgm:boot}

\begin{figure}[t]
  \centering
  \resizebox{\columnwidth}{!}{% F22: full lifecycle of a governed (ari_rqgm) run, set in the ARI house
% design language (_style.tex). Three phase columns, F21-style: BOOT
% (lavender / ideation fill) — mode resolution, the one-time founding
% registration transaction (25 prompts + 12 components), and the first
% epoch freeze; GOVERNED SEARCH (peach / execution fill) — the per-node
% hooks (proposal round, ungoverned node execution, governance-level
% assignment, adversarial round) inside a dashed loop enclosure; EPOCH
% BOUNDARY (pale green / gate fill) — audit, kernel-validated transition,
% frontier repair + clean room, and the next freeze. The dotted margin
% path at the bottom is the reflux into the next epoch's search. No
% accent edge: this figure has no single emphasis role (F23 carries it).
%
% style.tikzstyles is loaded by the preamble / the standalone wrapper;
% the ari* house styles come from _style.tex. The build cwd differs per
% entry point, so resolve _style.tex from each known prefix.
\InputIfFileExists{../shared/figures/tikz/_style.tex}{}{%
  \InputIfFileExists{shared/figures/tikz/_style.tex}{}{%
    \InputIfFileExists{report/shared/figures/tikz/_style.tex}{}{%
      \InputIfFileExists{_style.tex}{}{%
        \errmessage{F22: cannot locate figures/tikz/_style.tex}}}}}
\begin{tikzpicture}[aricanvas]

  % --- column 1: boot (mode + founding registration) -------------------
  \node[ariboxwide, fill=phaseidea]           (mode)     {\figlabel{F22.mode}};
  \node[ariboxwide, fill=phaseidea, below=of mode] (founding) {\figlabel{F22.founding}};
  \node[ariboxwide, fill=phaseidea, below=of founding] (freezeo) {\figlabel{F22.freeze0}};
  \node[ariphase, above=0.4cm of mode]        (h1)       {\figlabel{F22.head.boot}};
  \pic at ([xshift=-2.5mm]h1.west) {ariglyph bulb};
  \draw[ariarrow] (mode)     -- (founding);
  \draw[ariarrow] (founding) -- (freezeo);

  % --- column 2: governed search (per-node hooks, loop-enclosed) -------
  \node[aribox exp, right=of mode]            (proposal) {\figlabel{F22.proposal}};
  \node[aribox exp, below=of proposal]        (execute)  {\figlabel{F22.execute}};
  \node[aribox exp, below=of execute]         (level)    {\figlabel{F22.level}};
  \node[aribox exp, below=of level]           (advers)   {\figlabel{F22.adversarial}};
  \node[ariloop, fit=(proposal)(execute)(level)(advers)] (loopgov) {};
  \node[ariphase, above=0.65cm of proposal]   (h2)       {\figlabel{F22.head.search}};
  \pic at ([xshift=-2.5mm]h2.west) {ariglyph flask};
  \node[arinote, anchor=south west] at ([xshift=6mm]loopgov.north) (ttl) {\figlabel{F22.loop}};
  \pic at ([xshift=-2.5mm]ttl.west) {ariglyph loop};
  \draw[ariarrow] (proposal) -- (execute);
  \draw[ariarrow] (execute)  -- (level);
  \draw[ariarrow] (level)    -- (advers);

  % --- column 3: epoch boundary ----------------------------------------
  \node[aribox gate, right=of proposal]       (audit)    {\figlabel{F22.audit}};
  \node[aribox gate, below=of audit]          (transit)  {\figlabel{F22.transition}};
  \node[aribox gate, below=of transit]        (repair)   {\figlabel{F22.repair}};
  \node[aribox idea, below=of repair]         (freezen)  {\figlabel{F22.freeze1}};
  \node[ariphase, above=0.4cm of audit]       (h3)       {\figlabel{F22.head.boundary}};
  \pic at ([xshift=-2.5mm]h3.west) {ariglyph gate};
  \draw[ariarrow] (audit)   -- (transit);
  \draw[ariarrow] (transit) -- (repair);
  \draw[ariarrow] (repair)  -- (freezen);

  % --- handoffs between the phase columns -------------------------------
  \draw[ariarrow, rounded corners=5pt]
    (freezeo.east) -- ([xshift=8mm]freezeo.east) |- (proposal.west);
  \draw[ariarrow, rounded corners=5pt]
    (advers.east) -- ([xshift=8mm]advers.east)
    |- node[arinote, pos=0.3] {\figlabel{F22.note.trigger}} (audit.west);

  % --- reflux: the new epoch keeps searching ----------------------------
  \draw[arifeedback]
    (freezen.south) -- ([yshift=-7mm]freezen.south) -| (loopgov.south);
  \node[arinote] at ([yshift=-7mm]loopgov.south east) {\figlabel{F22.note.next}};

\end{tikzpicture}
}
  \caption{受治理运行的生命周期。启动阶段解析执行模式,执行一次性的
    创始注册事务(25 条提示词、12 个组件),并将创始在役集合冻结为
    纪元 000。纪元期间,逐节点钩子记录提案、指派治理级别(从仅固定
    检查的 L0 到全面裁决的 L3),并运行事件驱动的对抗轮;节点执行
    本身绝不会被否决。当纪元的节点计数触发器触发时,边界依次运行
    建议性审计、经内核校验的迁移,以及带净室再生成的前沿修复,随后
    冻结下一纪元的在役集合,搜索继续。}
  \label{fig:rqgm-lifecycle}
\end{figure}

受治理的运行分三步启动(\cref{fig:rqgm-lifecycle})。第一步,
\term{模式解析}:模式是一个双钥联锁——模式选择器与一个冗余的启用
标志必须一致,环境变量覆盖在赋值之前先行校验——任何不一致都会告警
并回退到无治理的默认,而不是去猜。解析出的模式及其来源(配置、环境
或恢复)记录于检查点。

第二步,运行时将搜索策略包裹进一个纯委托垫片:选择、扩展与剪枝的
行为与 \cref{sec:exploration} 完全一致,治理只是作为尽力而为的钩子
附着于其周围。

第三步——仅在全新的受治理检查点上——\term{创始注册}事务填充制度。
依据冻结在代码中的两张表,先是一个准备事件,继之以 25 条提示词注册
与 12 个组件注册,最后是单次提交。12 个创始组件是七个对抗者变体
(每种攻击类型一个,共享对抗者角色)、产物裁决者、辩护者、提案
路由器,以及两个元层智能体(提示词变异器与净室生成器);25 条创始
提示词还额外涵盖治理行动者(审计员、治理辩护人、治理裁决者)、
评审者模板,以及 \cref{sec:exploration} 的搜索模板,每条都连同其
模板字节的哈希一并注册。随后,该事务将创始在役集合——在役组件、
在役提示词哈希与效用策略——冻结为纪元 000 的不可变状态;其指纹
刻意不含时间戳,因而相同的制度会得到相同的指纹。创始是一次写入的:
恢复时,注册表通过重放事件日志重建,已提交的注册绝不重跑,而在
提交之前被中断的创始事务不可见,只会被简单地重新执行。

\subsection{逐节点治理}\label{sec:rqgm:node}

纪元期间,搜索与 \cref{sec:exploration} 中完全一致地运行;治理通过
三个尽力而为的钩子观察它。

\paragraph{提案。}
当扩展步骤生成候选方向时,提案路由器归档一份完整的\term{提案记录}
——本着「存下一切」的原则,包含完整的生成上下文——而搜索本身只
收到一份有上限的\term{提案摘要视图}。该摘要是搜索所见的\emph{唯一}
表面,因此提案生成器可以越变越丰富,而无需加宽通往搜索循环的接口;
归档的记录则始终可供审计与重放使用。

\paragraph{治理级别。}
每个被评估的节点都由一个纯触发阶梯指派一个治理级别——从 L0(仅
固定检查)到 L3(全面裁决)——依据的是确定性信号:跻身前沿最高分
之列、相对父节点的异常分数跃升、论文候选标志、新颖性主张,或评估器
置信度偏低。级别决定节点获得多少治理关注,而绝不决定它是否运行。
每次指派都追加到审计日志;当每纪元治理预算吃紧时,预算管理器可以
压低有效级别——降级的是治理,不是节点。

\paragraph{对抗轮。}
触发同一触发器族的节点还会经历一场\term{对抗轮}:七种对抗者类型
之一——夸大主张、证据缺口、指标操纵、可复现性、已有工作、成本爆炸
与提示注入——攻击该节点的\emph{产物}。对抗者从不攻击组件;攻击
记录在构造上就没有组件字段。原始攻击对分数是惰性的:只有在辩护者
作出回应、且产物裁决者维持该攻击之后,\term{已验证攻击记录}才存在;
也只有已验证的攻击才在纪元冻结的罚分策略下贡献一笔有界的效用罚分,
罚分前的分数保留于可加的指标键中。裁决者失败时默认攻击\emph{不成立}
——不罚分。已验证的攻击还会进入\term{对抗重放池}——之后考核候选
提示词的语料库(\cref{sec:rqgm:prompts})——这正是红皇后式的耦合:
今天被验证的攻击,就是明天的入学考试。对抗轮按节点幂等——案例日志
中的标记确保恢复后的运行绝不会重新攻击已处理过的节点。

\subsection{纪元边界}\label{sec:rqgm:boundary}

\begin{figure}[t]
  \centering
  \resizebox{\columnwidth}{!}{% F23: the epoch-boundary governance pipeline, set in the ARI house design
% language (_style.tex). Left column: the nine-step advisory audit
% (observe .. self-audit, numbered badges; the dashed loop enclosure marks
% it read-only w.r.t. registries). Right column: the report, the
% transition engine's resolution, the constitutional kernel's validation
% of every edge against the fixed T1-T21 table, the four-event boundary
% transaction, and frontier repair. The single accent edge (the rationed
% emphasis role of this figure) is the kernel's fail-closed block of an
% illegal transition.
%
% style.tikzstyles is loaded by the preamble / the standalone wrapper;
% the ari* house styles come from _style.tex. The build cwd differs per
% entry point, so resolve _style.tex from each known prefix.
\InputIfFileExists{../shared/figures/tikz/_style.tex}{}{%
  \InputIfFileExists{shared/figures/tikz/_style.tex}{}{%
    \InputIfFileExists{report/shared/figures/tikz/_style.tex}{}{%
      \InputIfFileExists{_style.tex}{}{%
        \errmessage{F23: cannot locate figures/tikz/_style.tex}}}}}
\begin{tikzpicture}[aricanvas]

  % --- column 1: the nine-step epoch audit (advisory) -------------------
  \node[aribox exp]                        (observe)  {\figlabel{F23.observe}};
  \node[aribox exp, below=of observe]      (reliab)   {\figlabel{F23.reliability}};
  \node[aribox exp, below=of reliab]       (evidence) {\figlabel{F23.evidence}};
  \node[aribox exp, below=of evidence]     (prosec)   {\figlabel{F23.prosecute}};
  \node[aribox exp, below=of prosec]       (defense)  {\figlabel{F23.defense}};
  \node[aribox exp, below=of defense]      (adjud)    {\figlabel{F23.adjudicate}};
  \node[aribox exp, below=of adjud]        (pool)     {\figlabel{F23.replaypool}};
  \node[aribox exp, below=of pool]         (selfaud)  {\figlabel{F23.selfaudit}};
  \node[ariloop, fit=(observe)(reliab)(evidence)(prosec)(defense)(adjud)(pool)(selfaud)] (loopaud) {};
  \node[ariphase, above=0.65cm of observe] (h1)       {\figlabel{F23.head.audit}};
  \pic at ([xshift=-2.5mm]h1.west) {ariglyph flask};
  \node[arinote, anchor=north] at ([yshift=-2.5mm]loopaud.south) (ttl) {\figlabel{F23.note.advisory}};
  \draw[ariarrow] (observe)  -- (reliab);
  \draw[ariarrow] (reliab)   -- (evidence);
  \draw[ariarrow] (evidence) -- (prosec);
  \draw[ariarrow] (prosec)   -- (defense);
  \draw[ariarrow] (defense)  -- (adjud);
  \draw[ariarrow] (adjud)    -- (pool);
  \draw[ariarrow] (pool)     -- (selfaud);

  % Step badges 1..8 (the report itself is step 9).
  \node[aribadge] at (observe.north west)  {1};
  \node[aribadge] at (reliab.north west)   {2};
  \node[aribadge] at (evidence.north west) {3};
  \node[aribadge] at (prosec.north west)   {4};
  \node[aribadge] at (defense.north west)  {5};
  \node[aribadge] at (adjud.north west)    {6};
  \node[aribadge] at (pool.north west)     {7};
  \node[aribadge] at (selfaud.north west)  {8};

  % --- column 2: report -> transition -> kernel -> commit -> repair -----
  \node[aribox paper, right=2.4cm of observe] (report) {\figlabel{F23.report}};
  \node[aribox idea,  below=of report]        (resolve){\figlabel{F23.resolve}};
  \node[aribox gate,  below=of resolve]       (kernel) {\figlabel{F23.kernel}};
  \node[ariboxwide, fill=phasepaper, below=of kernel] (txn) {\figlabel{F23.txn}};
  \node[aribox exp,   below=of txn]           (repair) {\figlabel{F23.repair}};
  \node[arinote, below=0.2cm of repair]       (retn)   {\figlabel{F23.note.retire}};
  \node[ariphase, above=0.4cm of report]      (h2)     {\figlabel{F23.head.transition}};
  \pic at ([xshift=-2.5mm]h2.west) {ariglyph gate};
  \node[aribadge] at (report.north west) {9};
  \draw[ariarrow] (report)  -- (resolve);
  \draw[ariarrow] (resolve) -- (kernel);
  \draw[ariarrow] (kernel)  -- (txn);
  \draw[ariarrow] (txn)     -- (repair);

  % --- the audit hands its report across the column gap -----------------
  \draw[ariarrow, rounded corners=5pt]
    (selfaud.east) -- ([xshift=10mm]selfaud.east) |- (report.west);

  % --- the one accent role: an illegal edge is blocked fail-closed ------
  \node[ariboxwide, right=2.2cm of kernel] (blocked) {\figlabel{F23.blocked}};
  \draw[ariarrow accent] (kernel) -- (blocked);

\end{tikzpicture}
}
  \caption{纪元边界流水线。九步审计是建议性的,对注册表只读;其报告
    输入注册表迁移引擎,后者解析出的迁移在四事件边界事务提交之前,
    由宪制内核对照固定的 T1--T19 表逐边校验。非法边以故障关闭方式
    被阻断——现任在役集合继续服务,运行继续。前沿修复仅在已提交的
    迁移使某个组件或提示词退役时运行。}
  \label{fig:rqgm-boundary}
\end{figure}

纪元冻结制度;边界是制度唯一可以变更的窗口。边界触发器是节点计数
(默认每纪元十个新节点),在每次外层搜索迭代的开头于主线程上检查
——绝不在有节点在途时检查——并在运行结束时再检查一次,因此最后
一个节点恰在最后一次迭代到达的纪元仍会事务性地关闭;未达触发数的
尾部纪元则简单地保持开放。运行结束时的这次检查之所以存在,是因为
实况运行暴露了这个缺口:恰好停在最后一个节点上的边界否则永远不会
触发,使该纪元的审计无从写下。

\paragraph{审计。}
关闭中的纪元首先接受审计(\cref{fig:rqgm-boundary}),由一条九步
流水线执行,其每个消耗 LLM 的步骤都有完全确定性的兜底,因此没有
模型的审计仍能走完。先从该纪元的记录中收集观察;可靠性评估为每个
受治理组件打分;\term{证据书记员}——且只有书记员——汇编证据包,
其中排除同角色指控,后者依宪制不可受理;继之为起诉,只有
\term{审计员}可以提交弹劾动议,并对照每纪元起诉预算缴存保证金,
动议获维持时保证金退还——这是对滥诉的经济制动;每份动议都会生成
一份辩护;裁决合议庭作出裁定;重放池得到更新;\term{治理自审计}
再经内核重新校验治理行动者自身的记录。结果是一份\term{治理报告}
——可靠性、动议、辩护、裁定、候选评估、建议、保证金账目——追加
到哈希链式审计日志。该报告是下一阶段的建议性输入;产出它的编排器
不能写注册表。

\paragraph{迁移。}
注册表迁移引擎将报告的建议与候选评估解析为一次纪元迁移——采纳、
处分、退役——它是冻结纪元状态、报告、注册表与所配置阈值的纯函数:
重跑时逐字节相同。组件与提示词沿一张十九条规则的固定表 T1--T19,
在十种状态(候选、已验证、影子、试用在役、在役、警告、察看、隔离、
退役、封禁)之间移动。内核校验解析出的迁移的每一条边:表外之边即
具名违规 \texttt{CK-REG-001};未具紧急形态而在纪元中途盖章的仅限
边界之边是 \texttt{CK-REG-002};并非由迁移引擎签发的迁移——比如
裁决者去写注册表——是 \texttt{CK-ROL-901}。被阻断的迁移以故障关闭
(fail-closed)方式中止,现任在役集合继续服务;被暂停的是治理,
运行永远不会。有效的迁移经由四事件边界事务提交——准备、关闭旧
纪元、开启新纪元、提交——因此任意两个事件之间的崩溃只留下一个
前缀,恢复时将其丢弃并确定性地重跑;同一迁移的重复提交则是受保护
的空操作。仅限边界变更的唯一例外是\term{紧急隔离}(规则 T16):
被查获内核级严重违规的在役组件可以在纪元中途被隔离,该紧急事件与
任何其他迁移一样被记录并校验。

\paragraph{前沿修复。}
如果已提交的迁移使某个组件或提示词退役,那么运行的过去便包含由
一个如今被判定有缺陷的制度所产出的记录,前沿不能继续为它们支付
利息。修复引擎计算被退役提示词哈希的依赖闭包——每一条实质派生自
它的记录——并通过可加标志(陈旧、可入前沿,以及擦除原因)把受
影响的节点标记为陈旧,经树文件持久化。擦除\emph{仅是逻辑性的}:
没有任何东西被物理删除,因此历史仍可审计,而前沿评分将其忽略。
修复是角色感知的:仅仅\emph{消费}了被退役评审者分数的效用,会在
\emph{原}纪元的冻结权重下用幸存输入重新计算;而被退役角色自己产出
的记录被径直作废——当原始权重无法找回时,节点被作废而非重新缩放,
因此任何数字都不会在另一套策略下被悄悄重新导出。如果修复自身未
通过校验,引擎保守降级,最终降至只收尾的排空模式:运行完成手头
未竟的工作,但不再扩展——绝不崩溃。

\subsection{演化中的评分}\label{sec:rqgm:utility}

此前每个小节都让制度演化,却将一样东西视为固定:\emph{评分}本身
——为前沿排序的效用策略,由一组坐标轴基底及其权重、一条合成规则、
一条前沿评分规则、一项深度罚项与一个探索常数构成。将评分在整个运行
生命周期内冻结是保守之选,但这并非红皇后一脉工作 \cite{iacob2026red}
真正作出的选择:其定义性主张是一种在每个纪元边界重写整个评分的树
搜索。受治理的效用演化弥合了这一缺口。效用策略本身成为一个自成一体
的可演化角色——与 \cref{sec:rqgm:prompts} 的提示词角色唯一的区别,
在于其现任是一份策略文档,而非提示词模板。效用策略在纪元内被冻结,
故一个纪元运行期间前沿由一条固定规则排序,且仅在边界处,由采纳提示
词的同一台迁移引擎重写。

一个专用提议器铸造候选策略,除此别无其他。它的四种默认变异是对边界
那份业已抽象的证据所作的纯算术——无模型、无时钟、无随机——因此默认
重写完全可复现,而且它被禁止读取前沿当前的分数:一个被调校得去讨好
自己已产出节点的评分,正是协同演化设计要防止的自指式合谋。提议器不
持有任何注册表权限,且其自身是一个被注册、可受处分的组件——它凌驾
不了自己所评分的角色,同样也没有任何绝对统治者凌驾于评分之上。

采纳搭乘迁移表中一条新边。一个已闯过验证与影子服务、且在冻结的重放
盘上得分不逊于现任的后继策略会将其\emph{取代}:规则 T20,唯一的
active 到 retired 之边,也是唯一不经隔离过渡的退役。内核仅对效用策略
角色允许它——每个行为型组件都保持保守的、仅处分的置换模型,对它们
而言直接退役仍被拒绝。取代将健康的现任以其\emph{旧的}策略哈希退役,
而由于如今每个节点都带着为它评分的那套策略的哈希,前沿修复
(\cref{sec:rqgm:boundary})便作废在退役策略下被排序的每一个节点,而
不仅是碰巧被罚分的那一子集。添加该边重新钉住了宪制哈希,一如修正案
理应如此(\cref{sec:rqgm:constitution})。

有一个限度值得像机制本身一样坦白地陈述。当评分的坐标轴是动态导出
的、且没有钉住固定的逐轴基底时,“得分不逊于”这道闸没有可供比对的
稳定次序,故取代退化为合法性与非退化性,而非一次经证实的改进。一道
能证明后继\emph{严格}更优的闸,需要一个锚定的基底——用以评分的真值
——那是 \cref{sec:rqgm:paper} 的论文流水线所供给的,而缺乏任何关于
研究质量的外部先知的探索阶段并不供给。这正是坦诚的红皇后姿态:没有
真值,准则便是在合法策略的空间内滚动,而非攀爬一座固定的山丘,我们
也不把它说成超出实情的东西。

\subsection{将攻击绑定到一个可问责组件}\label{sec:rqgm:targetbinding}

\cref{sec:rqgm:node} 的对抗轮在裁定之后产出一条已验证攻击记录。这条
记录一向指名它所牵涉的\emph{角色}——角色是攻击唯一能诚实获知的东西
——但至今它未曾指名一个可问责的\emph{组件},于是本应从一条已验证
攻击、经一次可靠性计数、走向一份弹劾动议的链条,无处可指而处于惰性。
缘由是结构性的:七个探索对抗者攻击的产物,由不持有已注册现任的节点
生成角色所作,故连原则上都不存在可以问责的组件。

如今由裁决者签发的记录,只要可解析,便带上单一的、注册表可解析的
目标:作出被该攻击所作废的决定的那个现任,在记录被铸造之时对照纪元
\emph{冻结}的在役集合解析而得——绝非一次在线注册表读取,那在一次
边界采纳之后会为前任的缺陷去弹劾后继。角色分离被严格保持:对抗者
攻击产物、且只攻击产物,而使一项观察成为可问责的,是裁决者的裁断
而非那项指控,因此没有谁既指控又绑定;一条会绑定其自身裁决者的记录,
会丢弃该绑定并三缄其口,而非抛出异常。当没有目标可解析时,该字段
缺席,记录与此前逐字节相同,因此探索对抗者的记录不变。

这一绑定在今日的生产环境中恰为一个对抗者而发火:\cref{sec:rqgm:paper}
的论文自我偏好攻击,其目标——稿件评审者——是一个已注册的创始组件。
七个探索对抗者依设计仍解析不出任何目标,因为它们攻击其产物的那个
角色不持有已注册现任;当这样一个角色某日获得现任,其绑定便无需任何
其他改动而随之成立,因为该机制是由角色表、而非任何逐对抗者的特例所
驱动的。

\subsection{提示词演化、净室与元层}\label{sec:rqgm:prompts}

\begin{figure}[t]
  \centering
  \resizebox{\columnwidth}{!}{% F24: the governed prompt lifecycle as a state machine, set in the ARI
% house design language (_style.tex). Column 1 (lavender / ideation fill)
% is the evolution side: the meta-tier clean-room generator and prompt
% mutator emit candidates, which climb the six-stage validation ladder
% (candidate -> validated -> shadow). Column 2 (peach / execution fill)
% is service: probationary_active -> active with the warning/probation
% sanction ladder below and the rehabilitation back-edge. Column 3
% (white) is the terminal sanction chain quarantine -> retired -> banned.
% The single accent edge (the rationed emphasis role of this figure) is
% T16 emergency quarantine, which forces an immediate close/open boundary.
% Dotted margin paths: exoneration re-entry (T18) and the
% clean-room regeneration loop from a retirement back to a fresh
% candidate. Rejection edges (candidate/shadow -> retired, T2/T5) are
% omitted from the geometry and stated in the caption.
%
% style.tikzstyles is loaded by the preamble / the standalone wrapper;
% the ari* house styles come from _style.tex. The build cwd differs per
% entry point, so resolve _style.tex from each known prefix.
\InputIfFileExists{../shared/figures/tikz/_style.tex}{}{%
  \InputIfFileExists{shared/figures/tikz/_style.tex}{}{%
    \InputIfFileExists{report/shared/figures/tikz/_style.tex}{}{%
      \InputIfFileExists{_style.tex}{}{%
        \errmessage{F24: cannot locate figures/tikz/_style.tex}}}}}
\begin{tikzpicture}[aricanvas]

  % --- column 1: evolution (meta tier emits, candidates climb) ----------
  \node[aribox idea]                       (cleanroom) {\figlabel{F24.cleanroom}};
  \node[aribox idea, below=of cleanroom]   (cand)      {\figlabel{F24.candidate}};
  \node[aribox idea, below=of cand]        (valid)     {\figlabel{F24.validated}};
  \node[aribox idea, below=of valid]       (shadow)    {\figlabel{F24.shadow}};
  \draw[ariarrow] (cleanroom) -- node[arinote] {\figlabel{F24.note.emit}} (cand);
  \draw[ariarrow] (cand)  -- (valid);
  \draw[ariarrow] (valid) -- (shadow);

  % --- column 2: service (adoption at the boundary; sanction ladder) ----
  \node[aribox exp, right=of cleanroom]    (probact)   {\figlabel{F24.probact}};
  \node[aribox exp, below=of probact]      (active)    {\figlabel{F24.active}};
  \node[aribox exp, below=of active]       (warning)   {\figlabel{F24.warning}};
  \node[aribox exp, below=of warning]      (probation) {\figlabel{F24.probation}};
  \draw[ariarrow] (probact) -- (active);
  \draw[ariarrow] (active)  -- (warning);
  \draw[ariarrow] (warning) -- (probation);
  \draw[ariarrow muted] (probation.west)
    to[bend left=28] node[arinote, pos=0.3] {\figlabel{F24.note.rehab}}
    (active.west);

  % adoption handoff: shadow (col 1) -> probationary_active (col 2)
  \draw[ariarrow, rounded corners=5pt]
    (shadow.east) -- ([xshift=8mm]shadow.east)
    |- node[arinote, pos=0.28] {\figlabel{F24.note.adopt}} (probact.west);

  % --- column 3: terminal sanctions --------------------------------------
  \node[aribox, right=of warning]          (quarantine) {\figlabel{F24.quarantine}};
  \node[aribox, below=of quarantine]       (retired)    {\figlabel{F24.retired}};
  \node[aribox, below=of retired]          (banned)     {\figlabel{F24.banned}};
  \draw[ariarrow] (quarantine) -- (retired);
  \draw[ariarrow] (retired)    -- (banned);

  % the one accent role: the forced emergency boundary edge (T16)
  \draw[ariarrow accent] (warning)
    -- node[arinote] {\figlabel{F24.note.emergency}} (quarantine);

  % exoneration re-entry (T18): back under probation
  \draw[arifeedback] (quarantine.north)
    |- node[arinote, pos=0.75] {\figlabel{F24.note.exon}} (probact.east);

  % clean-room regeneration: a retirement seeds a fresh candidate
  \draw[arifeedback]
    (retired.east) -- ([xshift=8mm]retired.east)
    |- ([yshift=7mm]cleanroom.north) -- (cleanroom.north);
  \node[arinote, above=0.25cm of probact] (reentry) {\figlabel{F24.note.reenter}};

\end{tikzpicture}
}
  \caption{受治理的提示词生命周期。元层智能体(提示词变异器;发生
    退役后为净室生成器)只产出候选。候选攀爬六阶段验证阶梯(静态、
    宪制、schema 试运行、对抗池重放、锚定、影子),且只能由迁移
    引擎在纪元边界处采纳,以试用身份进入服务。处分沿警告与察看逐级
    下行;着重标示的边是唯一的纪元中途迁移,即紧急隔离(T16)。
    退役提示词的文本被封存;其角色由基线兜底服务,直至净室替代者
    作为全新候选重新进入生命周期。从候选与影子到退役的拒绝边
    (T2、T5)为清晰起见予以省略。}
  \label{fig:prompt-lifecycle}
\end{figure}

制度提示词经由 \cref{fig:prompt-lifecycle} 的生命周期演化。
\term{提示词变异器}产出候选模板,除此之外别无其他——它不持有任何
注册表权限,而且从候选到在役刻意不设任何边,因此模型的任何输出都
无法直接进入服务。候选必须闯过六个阶段:\emph{静态验证}(确定性的
良构性检查,并确保未就地改动现任者)、\emph{宪制验证}(内核对候选
记录的检查)、\emph{schema 试运行}(一次注入调用必须产生可按该角色
声明的输出 schema 解析的回复)、对照对抗重放池的\emph{重放评估}
——候选要在击败其前任的那些攻击上接受评分——在固定任务集上的
\emph{锚定评估}(平局时偏向现任者),最后是\emph{影子}服务:候选
与现任者在抽样流量上并肩运行,效果仅限观察。至此,迁移引擎方能在
边界处采纳它,且采纳后以试用身份入役,而非在役。

退役触发全系统最严苛的路径。退役提示词的文本被封存:注册表视图将
其置为存根,一个专用的访问守卫在能力层拒绝对它的访问——无处例外,
永不解除。替代者在\term{净室}中产出:固定层组装一个封闭的输入包
——已提交的目录与抽象化的失效摘要,绝不含注册表中的提示词文本——
内核在生成之前对其筛查(违规属 \texttt{CK-CLN} 族,包括入场时对照
退役文本的词元瓦片(word-shingle)相似度筛查),且生成是一次性的,
没有任何可能去查找被封存文本的携工具循环。再生成的候选从最底一级
重新进入生命周期;其间,退役角色的席位由内建基线代为服务,因此受
治理的角色永远不会空缺。

元层受到比其所修改层级更严的治理。元智能体运行在只读沙箱中,身后
是一个仅白名单的工具代理,其唯一写入面是一个合成的提交工具;它们
的能力标志默认拒绝,危险标志恒为假;内核以子集运算检查\term{权限
不扩张}——元智能体不能授予后继者超出自身持有的权限;元层的每个
输出都以候选状态进入生命周期。元层可以提议;它不能使之生效。元
步骤也既非仅由失败触发,亦不沉默:变异器在每个纪元边界——包括
平静的纪元——都在按角色与按纪元的预算上限之下产出候选,而每个
边界的元步骤即使一无所提,也会记录一条审计可见的结果摘要。

\subsection{宪制及其修正}\label{sec:rqgm:constitution}

\begin{table}[t]
  \centering
  \small
  \begin{tabularx}{\columnwidth}{@{}lX@{}}
    \toprule
    \textbf{族} & \textbf{守护对象} \\
    \midrule
    \texttt{CK-SCH} & 记录 schema:畸形的治理侧记录阻断;畸形的
      节点侧记录告警。 \\
    \texttt{CK-HSH} & 哈希完整性:记录、已注册提示词哈希、产物
      摘要。 \\
    \texttt{CK-ACC} & 能力:按角色与层级的默认拒绝权限矩阵。 \\
    \texttt{CK-ROL} & 角色分离;\texttt{CK-ROL-901} 命名了典型
      违规——除迁移引擎之外的任何人写注册表。 \\
    \texttt{CK-REG} & 对照 T1--T19 的迁移合法性:\texttt{001} 为
      表外之边,\texttt{002} 为在纪元中途尝试仅限边界之边。 \\
    \texttt{CK-EPO} & 纪元不变性:在役集合不得在纪元中途漂移。 \\
    \texttt{CK-ERA} & 选择性擦除:陈旧记录不得进入前沿评分。 \\
    \texttt{CK-AUD} & 审计日志完整性:事件哈希链及其单调性。 \\
    \texttt{CK-CLN} & 净室卫生:输入包筛查与对照封存文本的沾染
      检查。 \\
    \texttt{CK-CTX} & 上下文作用域:角色作用域的视图不得泄漏出其
      角色。 \\
    \bottomrule
  \end{tabularx}
  \caption{宪制内核的违规代码族。代码是稳定标识符,从不重新编号;
    每个校验器报告其触发的具体代码,由执法适配器按上下文决定告警
    还是阻断——状态变更阻断,节点侧发现告警。}
  \label{tab:ck-codes}
\end{table}

宪制是代码,不是配置。迁移表、能力矩阵、严重度与擦除规则,以及
净室筛查都存于源码之中,由这些表导出的单一\term{宪制哈希}被一个
守卫测试钉住。改动任何规则都会改变哈希并使守卫失败,因此宪制变更
不可能作为副作用发生:作者必须显式地重新钉住哈希,并附上注明日期
的理由——这是一次深思熟虑、可供评审的行为,与其说是配置变更,
不如说更接近修正案程序。此类修正迄今发生过一次:角色词汇表中存在
的一个元层角色在能力矩阵中缺失,导致对它的每次能力查询都被拒绝;
该修正案以与其权限最窄同侪相同的最小授权添加了该角色,并在理由
就位的情况下重新钉住了哈希。数值阈值(预算、触发计数)是可配置的;
表的拓扑不是。宪制的一份人类可读呈现被复制进每个受治理检查点,
其哈希记录于该次运行的元数据——修改该副本不会改变任何东西,这
正是设计使然。\Cref{tab:ck-codes} 列出内核报告的违规代码族。

\subsection{治理论文流水线}\label{sec:rqgm:paper}

治理迄今包裹的是探索搜索。同一套机制也治理 \cref{sec:paper} 的稿件
撰写流水线,置于一条与第一条完全独立的第二执行轴之下。论文阶段的
模式在两者间选择:\emph{线性}流水线——默认,与 \cref{sec:paper}
逐字节相同,论文路径上不构造任何治理对象——以及将论文撰写置于同一
套纪元约束治理之下的\emph{存档}模式。一如 \cref{sec:rqgm:boot} 的
探索模式,它是一个双钥联锁,模式选择器与一个冗余的启用标志必须一致,
而两条轴相互正交:治理与否的探索,同线性与否的存档式论文撰写,任意
组合皆为合法。

存档模式为论文撰写实现了这一方法赖以定义的四项承诺。\emph{其一,可
重写评分之下的树搜索。}该模式是一棵在草稿空间上货真价实的浅层最佳
优先树:第一层是一小撮不同立意的种子草稿,更深的节点是迭代式的精修,
并以受治理评审者的合成分数外加一项奖励欠采样立意的多样性加成,作出
真正的前沿选择。它原封不动地复用探索搜索的剪枝与计数——总节点预算
先于深度截断生效,故更深的树只是重新分配同一批节点,而非将其相乘
——评审者效用在受治理的效用演化(\cref{sec:rqgm:utility})之下跨纪元
协同演化,前沿修复则作废在退役策略下评分的草稿。\emph{其二,对抗者
攻击评价,而不止攻击产物。}第八种对抗者类型攻击评审者的过度接受:
一份现任评审者给了高分、而人工标注的锚定语料本会拒绝的 AI 撰写草稿。
锚定只供给前置信号——该攻击哪些被过度接受的草稿——而由一场货真价实
的对抗者—辩护者—裁决者之轮产出已验证攻击记录;径直从锚定去制造那条
记录,恰是设计所禁止的那种合谋。\emph{其三,没有绝对统治者。}两个
受治理角色,稿件撰写者与稿件评审者,驱动论文技能,而技能始终是不受
治理的「手」,绝不跨进程被治理;两个角色都被注册、可受处分、可演化,
且仅在存档模式下注册,因此探索的启动逐字节相同。评审者身为评价者,
以红皇后论文中评价者相同的方式赢得信任:与一份留存的、人工标注的
接受/拒绝锚定语料相一致。撰写者同样被锚定——锚定于红皇后论文的
撰写者从未拥有过的真值。由于 ARI 书写的是真正运行过的实验,第 0 层的
主张—证据硬闸确定性地为一份草稿的忠实度评分:把其执行接地的主张率、
数值主张的可复现率与数值覆盖率折叠为单一分数;一份对已核验发现作出
过度主张或误报的稿件,独立于任何演化中的评审者,客观上更差。因此
ARI 的论文撰写者是红皇后论文\emph{编码}域的形态——一个确定性验证器
加上一个协同演化的评审者——而非该论文中无锚的论文撰写域;人工精选
语料仍然只是评审者的锚定,故撰写者的锚定不需要任何额外数据成本。
\emph{其四,一切都受宪制约束。}两个角色中的任一个
只经由 \cref{sec:rqgm:targetbinding} 的目标绑定所闭合的、受处分的
可靠性—起诉—弹劾之路开启;被取代的论文提示词并非退役,而是经由
一条论文角色专属的 active 到 shadow 之边(规则 T21)降为可复任的
影子待命;而锚定对其真值中可被机器自举的比例之上限,由代码强制执行
——一旦逾越便拒绝该语料并降级,绝不抛出异常。

对撰写者的处分是一场货真价实的对抗之轮,绝非制造出来的记录。一份被
过度接受、且主张闸也判定为不忠实的草稿有两个有责组件——接受它的评审者
与产出它的撰写者——因此自我偏好之轮按每个可解析角色发出一条已验证攻击,
撰写者的绑定按草稿以该草稿的忠实度为闸;一份被过度接受但忠实的草稿
只绑定评审者。忠实度分数以撰写者的提示词哈希为键,而非以其组件身份为键,
因为一个组件身份会跨越接连的提示词版本,以组件为键的分数会把现任者的
失败泄漏给它的后继。闸本身只被读取:它以不持久化的模式被调用,绝不被
包裹、编辑或治理。

向结果核验的交接被刻意保持原样。存档中最优的草稿被一次性写入检查点
的正典稿件文件,而 \cref{sec:paper:contract} 那道既有的编译兼主张—
证据硬闸——一条固定的、绝不被内核包裹的第 0 层核验路径——在与线性
流水线所用完全相同的契约之下对它运行。治理选择哪一份草稿去面对这道
闸;它绝不放松这道闸。成本以与探索阶段相同的纪律保持有界:每纪元的
草稿群体由一份节点预算封顶,而该预算在任意深度都映射到搜索的总节点
截断,至于降级的入门路径——无锚定语料、评审者协同演化被抑制——则是
约以线性流水线的成本作 N 中择优的评审草稿。

四个限度圈定存档模式的主张范围。两个提示词都会协同演化,但撰写者的
置换在构造上是保守的:撰写者的后继会攀升至影子位并在那里等待,直到
现任者在主张闸的忠实度上\emph{退化}——因此,一个仅仅看起来更好的
挑战者永远不会置换一个忠实的现任者。这就是行为角色模型:开启一个角色
的是处分,而不是比较。另一方面,撰写者的\emph{草稿}赢家仍然是纪元
局部的:它们由纪元内冻结的评审者排序,故在不同评审者版本下评分的草稿
从不相互比较。针对撰写者的攻击搭乘自我偏好之轮,而该轮仅在存在被过度
接受的锚定案例时才发火,故一个在不过度接受任何东西的评审者之下工作的
不忠实撰写者,今日不会被处分;一个专用的撰写者对抗者是另一项决定。
锚定语料默认关闭,故默认的存档模式
在提供一份经甄选的接受/拒绝语料之前,就是 N 中择优的评审草稿——而且
由于撰写者的忠实度证据落在同一个锚定池上,没有语料也就意味着没有对
撰写者的处分,尽管撰写者自身的锚定并不需要任何自己的甄选数据。在每
个论文阶段两轮这一默认下,边界与弹劾机制会发火,但一次采纳
所需的多阶段攀爬——约五个边界——并不完成,故一次必须目睹在役提示词
发生变化的运行需要更多轮次。而在一轮之内将一份真实的被过度接受草稿
降级的相内罚项被推迟;今日发火的是问责与协同演化通道,而非
将该草稿就地降级。

\subsection{可观测性}\label{sec:rqgm:observability}

治理留下的审计踪迹与搜索本身同属一类:只追加的事件日志作为事实,
派生快照用于快速读取,全部位于检查点之内。模式来源记录缺失,即
意味着一次纯粹的无治理运行。模式来源记录(\texttt{rqgm\_state.json})
保存解析出的模式及其来源;注册表快照(\texttt{rqgm\_registry.json})
是当前的组件与提示词表,恢复时从事件重建;迁移日志
(\texttt{rqgm\_transitions.jsonl})承载各注册与纪元事务,冻结纪元
快照(\texttt{epoch\_state.json})保存在役集合、效用策略与无时间戳
指纹;审计日志(\texttt{rqgm\_audit.jsonl})是哈希链式的制度史——
每个事件都携带自身哈希与前驱哈希,截断或篡改都会破坏链条,恢复时
会重新验证;对抗案例日志(\texttt{rqgm\_adversarial\_cases.jsonl})
记录原始攻击、辩护、裁定、已验证攻击与逐节点幂等标记;提案存储
(一个 \texttt{proposals/} 目录)以索引归档完整的提案记录;演化出
的提示词正文位于 \texttt{rqgm\_prompts/} 之下,一次写入并经哈希
校验;提示词踪迹(\texttt{prompt\_trace.jsonl})为每条渲染出的
提示词打上所服务的提示词版本戳,创始模板服务期间为 null;擦除汇总
(\texttt{rqgm\_erasure\_state.json})是擦除事件的纯折叠——它不
存在即意味着没有任何东西陈旧。

三段来自实况检查点的摘录勾勒出这些记录的形状(经裁剪并轻度重排;
可识别的运行名已隐去)。创始注册,摘自迁移日志:

\begin{Verbatim}[breaklines, fontsize=\scriptsize]
{"event_type": "epoch_transaction_prepare",
 "payload": {"transition_id": "transition_founding"}}
{"event_type": "component_registered", "payload":
 {"component_id": "adversary_overclaim_v1",
  "role": "adversary", "tier": "institutional",
  "status": "active", ...}}
\end{Verbatim}

\noindent
一条逐节点治理级别指派,摘自哈希链式审计日志:

\begin{Verbatim}[breaklines, fontsize=\scriptsize]
{"event_type": "governance_level", "payload":
 {"epoch_id": "epoch_000", "node_id": "node_a37b4c06",
  "level": 3, "triggers": ["top_k"]},
 "event_hash": "ae8eeb7521f1",
 "prev_event_hash": "23567870882b"}
\end{Verbatim}

\noindent
以及来自评估工具链一次注入运行(\cref{sec:rqgm:empirical})的两次
检出:

\begin{Verbatim}[breaklines, fontsize=\scriptsize]
{"event_type": "constitutional_violation", "payload":
 {"check": "validate_record_schema",
  "codes": ["CK-SCH-N01"], ...}}
{"event_type": "prompt_candidate_rejected", "payload":
 {"stage": "schema_dry_run", "failure_count": 4, ...}}
\end{Verbatim}

\subsection{经验现状}\label{sec:rqgm:empirical}

与 \cref{chapter:eval} 一样,我们报告这套机制当下能佐证的内容,并
推迟分布性的结论。今天存在三类证据。第一,单元与集成测试套件——
撰写本文时共 3802 个测试,其中 750 余个位于专门的治理套件——直接
钉住各项宪制不变量:非法迁移被阻断、原始攻击对分数惰性、注册表
单一写入者、擦除陈旧性闭包、净室沾染筛查、效用策略取代之边仅对该
角色合法,以及探索与论文两个无治理默认的行为皆保持不变。

第二,实况受治理运行在真实检查点上端到端地演练了完整生命周期。在
一项小型 HPC 内核优化实验的探索运行中,创始事务提交了其 25 条提示词
与 12 个组件的注册;每个被评估的节点都获得了留档的治理级别;第一个
纪元经四事件边界事务关闭;边界审计产出了一份治理报告,其自审计覆盖
了四个治理行动者,发现内核违规为零。受治理的效用演化跨真实边界运行:
在役效用策略哈希从一个纪元变到下一个——一个被观察到的序列历经
\texttt{fed4460f44f6},继而 \texttt{8a9da4fb9dc3},继而
\texttt{8ea0dac3cd1b}——纪元指纹随之改变,健康的现任因取代而退役,
而在退役策略下评分的每个节点都对前沿被作废。一次论文存档运行驱动
两个受治理的论文角色协同演化:在役评审者提示词哈希经由一场货真价实的、
产出一条绑定到现任评审者的已验证攻击记录的对抗者—辩护者—裁决者之轮,
从 \texttt{04b3c49d070d} 移到 \texttt{c0dfd764abab},而存档中最优的
草稿抵达了真实的编译兼主张之闸。以一个发出不忠实草稿的撰写者驱动同一
运行,则使在役撰写者提示词哈希在八轮之内从 \texttt{f38a15f0f140} 移到
\texttt{b2c36f9a8232}——四轮平坦,随后采纳——经由一场其已验证攻击记录
点名现任撰写者为目标的货真价实之轮。两个对照证实了这些耦合是承重的
而非附带的。一旦扣下锚定语料,便不产生过度接受信号,不验证
任何攻击,也不发生任何采纳。而仅把主张闸的忠实度翻转到最大,则产生
零次撰写者攻击、零次撰写者动议、撰写者哈希恒定——与此同时评审者依然
被攻击并被弹劾,撰写者的后继依然停在影子位——因此撰写者之未被采纳,
是由处分的缺席所致,而非由一次死掉的运行所致。

第三,评估工具链的注入层——经真实流水线重放的确定性故障注入——演示
了检出能力:一个重复裁决的裁决者产生了三条归因于被注入组件的已验证
攻击记录,一条畸形的生成器记录触发了 schema 检查,一个退化的变异器
候选在 schema 试运行阶段被拒——这些都作为上文摘录的审计事件可见。

这些运行确立了每套机制都会发火;它们尚不构成治理改善研究的证据,且
其中若干是在其完整效应得以显现的阈值之下运行的。撰写者仅在退化时才被
采纳,而绝非因挑战者更优,且其处分搭乘评审者的过度接受之轮;在两个论文
轮次这一默认下,弹劾机制发火,
却未完成一次需要更多边界的采纳;论文锚定默认关闭,故默认的存档
模式是 N 中择优的评审草稿;而目标绑定链条在生产环境中仅对论文自我
偏好对抗者发火,探索对抗者的链条依设计保持惰性。尚不存在的是受控
比较——受治理与无治理的搜索质量在一个种子与实验面板上的对照;该研究
与 \cref{sec:eval:limits} 所推迟的评估同属一处。
